\section{Open Estimates And Their Proven Reductions}
\label{sec:open-estimates}

The filter-seven excess bound and copy-block excess control properties are the last two steps of a longer argument, and the
article's conclusion (Section~10, Section~12) identifies two open estimates at the live
twin-prime frontier. This section introduces them so that the article's
conclusion is readable on its own. It states each estimate's exact reduction,
open bound, and relationship to what the article has already proved. Further
algebraic reductions (activation shells, CRT lift indices, and Gram matrices)
exist beyond the article's scope. This section states the exact reductions
needed to follow the article's own arc and keeps the remaining estimate open.

\subsection{Accepted-Boundary Discrepancy Estimate}
\label{sec:accepted-boundary-discrepancy}

The reduction below is proved for every nonempty conditioned chain
$5\le r_0<\cdots<r_{m-1}<Q$, over accepted anchor values in one interval,
followed through the conditioned filter chain; the terminal
signed-mean-square estimate is open.

Recall from Section~5.1 that the harmful excess at layer $i$ is
$b_i=a_iN_i-N_{i+1}$. The accepted-boundary model isolates the \emph{strike-density error}
$\varepsilon_i=H_i/A_i-1/r_i$, where $H_i$ is the count of accepted anchors
struck by filter $r_i$ and $A_i$ is the accepted-anchor population. Its
central result is an exact bridge: the expected bulk density $1/r_i$ is
already exact, and the harmful excess reduces to a difference of signed
M\"obius boundary sums. Specifically, with
$E_i=A_i-\ell\varphi(P_i)/P_i$ the centered inclusion--exclusion discrepancy
at layer $i$,
\begin{equation*}
H_i-\frac{A_i}{r_i}
=
\left(1-\frac1{r_i}\right)E_i-E_{i+1},
\qquad
\varepsilon_i
=
\frac{(1-1/r_i)E_i-E_{i+1}}{A_i}.
\end{equation*}

The signed discrepancies telescope under one-anchor survival weights, but the
weighted-energy budget requires a \emph{weighted sum of their squares}:
\begin{equation*}
\sum_i
w_i\frac{r_i}{2(r_i-2)}
\left(\left(1-\frac1{r_i}\right)E_i-E_{i+1}\right)^2.
\end{equation*}

Bounding the divisor summands independently gives only
$|E_P|\lt2^{\omega(P)}-1$, exponentially too large. Success requires signed
cancellation, correlation between consecutive layers' boundary sums, or
cross-layer averaging --- a genuinely new arithmetic input. The properties from Strike Divisor-Activation Kernel through First-Deletion Reindexing
catalog the activation-shell, CRT-lift, summatory, Gram, and first-deletion
reductions; each returns the original energy after exact algebra, so the
remaining input is signed arithmetic, not another coordinate rewrite.

\textbf{Why this is the article's general coefficient.} The filter-7 calculation
of Section~7 is the one-layer, one-prime instance of this same discrepancy:
$|b_7|\le18/7$ came from exact residue \emph{order}, and the general layer
coefficient $b_i$ is exactly the two-residue boundary discrepancy studied
here. A 

\subsection{Residue-Collision Energy Estimate}
\label{sec:residue-collision-energy}

For every incoming prime $r\ge5$ and its actual conditioned layer population,
over 2-gap-start residues modulo one incoming prime in one conditioned layer,
the reduction below is proved; the relative four-point correlation estimate
is open.

The residue-collision energy is the input that the copy-block bridge (Section~9)
consumes. Let
$c_t$ count the incoming 2-gap starts in residue class $t\bmod r$, so
$N_r=\sum c_t$. The centered deviation is $d_t=c_t-N_r/r$, and the
residue-collision energy is
\begin{equation*}
V_r=\sum_{t\bmod r}d_t^2.
\end{equation*}

The two harmful classes contain $c_0+c_{-2}$ starts. Their centered excess
reduces exactly to the histogram second moment and its autocorrelation:
\begin{equation*}
C_r
=
\sum_{t\bmod r}c_t^2
=
N_r+2\sum_{h\ge1}A_r(6rh),
\end{equation*}

where $A_r(\cdot)$ is the four-point autocorrelation of the start indicator
at the given shift. The open estimate is the \emph{relative} bound
\begin{equation*}
C_r\le N_r+\frac{N_r^2}{r},
\qquad\text{equivalently}\qquad
V_r\le\frac{N_r^2}{r}.
\end{equation*}

An absolute upper-bound-sieve estimate is insufficient until its
normalization by the actual $N_r$ is justified independently. Minimal
falsifying histograms exist at small scale ($3+2+1$ at $(r,N)=(5,6)$;
$2+2$ at $(7,4)$), but exact conditioned-layer search through $Q\le251$
found none. A 

\textbf{How the two compose.} The copy-block bridge of Section~9 proves that the
complete-block harmful excess $B_j=d_{t_j}+d_{t_j-2}$ satisfies
$\sum_jB_j^2\le4V_r$. Therefore a relative bound for the residue-collision
energy $V_r$ immediately controls the complete-block portion of the terminal
harmful excess. The frontier is consequently a single composition: a relative
residue-energy estimate feeding the signed boundary discrepancy through the
copy-block bridge, with the two partial old-period boundary fragments still
open. This is exactly what the article's conclusion (Section~10) names.
